\chapter*{Introduzione}
\addcontentsline{toc}{chapter}{Introduzione}
\label{chap:introduzione}

\section*{Inquadramento del problema}
% Linee guida: L'ottimizzazione delle risorse e della logistica temporale nelle organizzazioni complesse.
% Inserisci qui l'introduzione al tema dello scheduling e dell'ottimizzazione in contesti organizzativi.

\section*{Il problema dello School Timetabling}
% Linee guida: Definizione generale dell'istituzione scolastica come dominio applicativo di problemi di scheduling combinatorio.
% Descrivi qui le peculiarità del contesto scolastico e perché rappresenta un problema di ottimizzazione combinatoria.

\section*{Obiettivi della tesi}
% Linee guida: Modellizzazione formale di un caso di studio reale mediante strutture discrete e analisi comparativa delle prestazioni algoritmiche.
% Descrivi gli obiettivi specifici del lavoro svolto sul caso reale dell'Istituto Saffi-Alberti di Forlì.

\section*{Struttura della Tesi}
% Descrizione sintetica della ripartizione dei capitoli successivi.
